\documentclass[11pt]{article}
\usepackage{ochamath}
\subjectcolor{2F5DA8}
\setsheet{線形代数}{基底と次元　演習}{https://math.ochanote.com/linear-algebra/basis-dimension/}
\namelinetrue
\begin{document}
\makesheethead{問題}

\problem[基底座標]
$\boldsymbol p_1=(1,1)^{\mathsf T},\boldsymbol p_2=(1,-1)^{\mathsf T}$ による基底で、$\boldsymbol x=(6,2)^{\mathsf T}$ の座標を求めよ。
\workspace{38mm}

\problem[列空間と零空間]
$A=\begin{pmatrix}1&2&0\\0&0&1\\1&2&1\end{pmatrix}$ の列空間・零空間の基底をそれぞれ求めよ。
\workspace{38mm}
\newpage

\problem[基底の判定]
$W=\{(x,y,z)^{\mathsf T}:z=x+y\}$ に対して $(1,0,1)^{\mathsf T},(0,1,1)^{\mathsf T}$ は基底か。次元も答えよ。
\workspace{38mm}

\problem[測定の自由度]
4成分の入力を持つ行列 $A$ の階数が3である。出力を変えない入力方向の空間の次元を求め、一意な復元が可能か述べよ。
\workspace{38mm}
\end{document}
